% 19 November 2026. Prepared for Edward Baker with Codex assistance.
% Lead gpt-6-astra/max. Existing comparisons are rederived, not new claims.
\section{What a backward estimate must control}
\label{sec:continuation}

\subsection{The cumulative drift budget}

The pair law of Section~\ref{sec:pair-chain} reads
\begin{equation}
 D_t=2-4DM,\qquad
 D_u=-\mathcal E,\qquad
 \mathcal E=2-4DM=2-g^2M.
 \label{cont:drift}
\end{equation}
The target is a time integral of the actual coupled quantity, not a
sign inferred from the approximation $d^2M\approx2$. Even
$g^2M=2+O(\varepsilon)$ does not select a sign for $D_u$; it selects
an error scale whose accumulated effect must be compared with the
initial squared gap.

\begin{proposition}[Cumulative continuation margin]
\label{cont:cumulative}
Suppose a counted pair satisfies \eqref{cont:drift} on a stopped path
$0\le u<\tau$, $D>0$, and $D(0)\ge d_0>0$. If an integrable real
function $\varepsilon$ satisfies $\mathcal E(u)\le\varepsilon(u)$,
then
\begin{equation}
 D(u)\ge d_0-\int_0^u\varepsilon(v)\dd v.
 \label{cont:integral}
\end{equation}
If the supplied inputs persist up to a possible first target merger
and
\begin{equation}
 \sup_{0\le s\le\Delta}\int_0^s\varepsilon(v)\dd v<d_0,
 \label{cont:strict}
\end{equation}
that merger cannot occur by $\Delta$. For a nonnegative allowance it
suffices to check the integral at $\Delta$.
\end{proposition}

\begin{proof}
Integrate \eqref{cont:drift}. The proposed first merger has $D\to0$,
whereas \eqref{cont:strict} leaves a uniform positive lower bound.
No differential equation is divided by $D$ at the event.
\end{proof}

This elementary comparison specifies a quantitative task. It does not
supply $\varepsilon$, a root-count tube, or coverage of other gaps.
When the allowance is signed, checking only its final cumulative value
could miss an earlier exhaustion; the supremum in
\eqref{cont:strict} is essential. At height $X$, with
$L=\log(X/(4\pi))$ and $D(0)$ of order $L^{-2}$, even a small
absolute drift error must be integrated on that scale.

\subsection{What initial interaction alone buys}

In an all-real simple configuration, exterior terms for an adjacent
pair are positive. Put
\begin{equation}
 \mathcal F=12S_2+16DS_3.
 \label{cont:F}
\end{equation}
The first moment equation gives
$M_u=-2M^2+\mathcal F\ge-2M^2$.
If $b_0=M(0)\ge0$, comparison with the scalar Riccati solution gives
$M(u)\ge b_0/(1+2b_0u)$. Retaining this lower bound in
$D_u=-2+4DM$ and using the factor $(1+2b_0u)^{-2}$ yields
\begin{equation}
 D(u)\ge(1+2b_0u)
       \{D(0)-2[1-b_0D(0)]u\}.
 \label{cont:initial-mass}
\end{equation}
The comparison is valid before its all-real stopped regime ends.
When $b_0D(0)<1$, its strict positive duration is
\begin{equation}
 u<\frac{D(0)}{2[1-b_0D(0)]}.
 \label{cont:initial-horizon}
\end{equation}
This derivation also covers $b_0=0$ by continuity. For an initial
near-clock configuration with $b_0D(0)\to1/2$, the duration approaches
$D(0)=g(0)^2/4$. It still vanishes at high height. This is the allowance
of this particular comparison, not a dynamical prediction of collision
\cite[\S\S13--15]{VCDCMManuscript}.

For example, an exact cosine heat field has stationary equally spaced
zeros, $M=2/g^2$ and $D=g^2/4$. It survives indefinitely, while
\eqref{cont:initial-horizon} stops at $g^2/4$. The forcing discarded
in the Riccati comparison balances $2M^2$ exactly. This comparison
control has uniform density; the canonical reference in
Section~\ref{sec:pair-chain} remains logarithmically height dependent.

\subsection{Keeping the joint forcing}

Define
\begin{equation}
 V=1-2DM,\qquad K=2D^2\mathcal F
                      =24D^2S_2+32D^3S_3.
 \label{cont:VK}
\end{equation}
These quantities combine gap size and exterior moments before an
inequality discards their dependence. Direct differentiation gives
\begin{equation}
 D_u=-2V,\qquad V_u=\frac{1-V^2-K}{D}.
 \label{cont:VK-law}
\end{equation}
Indeed $V_u=4M-4DM^2-2D\mathcal F$, and substitution of
$2DM=1-V$ gives the second equation. This is a regular calculation
for $D>0$; the comparison below is stopped before a merger.

\begin{theorem}[Joint forcing barrier]
\label{cont:barrier}
Suppose the real complete pair laws hold on a stopped path with $D>0$,
$V(0)\le v_0$, and $K(u)\ge1-\eta^2$ for a constant $\eta\ge0$.
For $v\ge\max\{v_0,\eta\}$,
\begin{equation}
 V(u)\le v,\qquad D(u)\ge D(0)-2vu.
 \label{cont:barrier-bound}
\end{equation}
A positive initial lower bound $d_0$ protects the target pair through
$\Delta$ if $2v\Delta<d_0$ and the stipulated inputs persist up to
a first event.
\end{theorem}

\begin{proof}
On each compact pre-event interval the coefficient below is continuous.
For $w=V-v$, \eqref{cont:VK-law} implies
\[
 w_u+\frac{V+v}{D}w
 \le\frac{\eta^2-v^2}{D}\le0.
\]
A positive integrating factor and $w(0)\le0$ prove $w\le0$.
Integrate $D_u=-2V$, then exhaust the stopped path by compact
subintervals. A first merger contradicts the resulting strictly positive
lower bound as in Proposition~\ref{cont:cumulative}.
\end{proof}

The proof itself needs no positive summands: a signed complete bound
for $K$ in a conjugation-invariant counted tube also suffices. Positive
omission of a tail, however, requires the all-real adjacent-pair regime.
If another pair collides first and ends that regime, it ends that use of
positivity. Protection of one pair is not protection of the complete
population.

For $D(0)\asymp L^{-2}$, sufficient rates to pay a fixed duration are
\begin{equation}
 V(0)_+=o(L^{-2}),\qquad
 \sup_u(1-K(u))_+=o(L^{-4}).
 \label{cont:rates}
\end{equation}
The square root of the forcing deficiency explains the stronger second
rate. These are sufficient requirements of this barrier, not necessary
conditions for actual noncollision.

\subsection{Finite heads: an exact comparison and its perturbation}

For an equally spaced comparison configuration with target endpoints
$0,g$, the $m$th exterior occurrence on either side has
$Q_m=m(m+1)g^2$. With $D=g^2/4$ and $J$ such occurrences on each side,
let $M_J,K_J$ denote their contributions. Telescoping gives
\begin{equation}
 2DM_J=\frac{J}{J+1},\qquad
 K_J=1-\frac1{(J+1)^3}.
 \label{cont:lattice-head}
\end{equation}
For the second identity the contribution from both sides is
\[
 \frac3{m^2(m+1)^2}+\frac1{m^3(m+1)^3}
   =\frac1{m^3}-\frac1{(m+1)^3}.
\]
These exact formulas are local comparison tools, not a uniform-density
model of the complete zeta divisor.

They have a useful robust version. In the all-real adjacent-pair stopped
regime, suppose the actual target gap obeys
$g\ge(1-e)d$, $0\le e<1/3$, and the $m$th selected exterior zero
on each side is at distance $r$ from the nearer endpoint, with
\begin{equation}
 r\le(1+2e)md,\qquad r+g\le(1+2e)(m+1)d,
 \quad 1\le m\le J.
 \label{cont:head-input}
\end{equation}
Writing $\gamma=(1-e)/(1+2e)$ gives
$D/[r(r+g)]\ge\gamma^2/[4m(m+1)]$.
Since the quadratic and cubic powers in $K$ are nonnegative,
\begin{equation}
 K\ge1-18e-(J+1)^{-3},\qquad
 V\le6e+(J+1)^{-1}.
 \label{cont:head-bounds}
\end{equation}
To check the constants, use $1-\gamma\le3e$,
$1-\gamma^6\le18e$ and $1-\gamma^2\le6e$ in
\eqref{cont:lattice-head}, and omit only positive exterior terms.
A uniform version of the first bound and the initial instance of the
second supply Theorem~\ref{cont:barrier}. Adding more initial neighbors
alone does not provide the whole-time first bound.

\subsection{What the existing canonical acquisition already uses}

The source analysis in \cite[\S15]{VCDCMManuscript}, based on
\cite[Theorem 1.5 and \S9]{Polymath2019}, works on a fixed slab
$0<t_0\le t\le t_*\le1/2$. The height onset depends exponentially
on $1/t_0$. The positive-head acquisition is restricted to the globally
all-real simple stopped portion, begun from an all-real simple configuration. On that same slab it first locates individual consecutive
zeros near the source clock; it then estimates complete $M$ and a
neighboring head for $K$. With $L=\log(X/(4\pi))$, a head of
order $L^3$ neighbors and its tail payment give
\[
 |V|=O(L^{-3}),\qquad (1-K)_+=O(L^{-9})+o(L^{-m})
\]
for each fixed $m$, after the sufficiently large height is chosen.
The error terms inherited from localization are smaller than every
fixed inverse power of $L$ on this fixed positive slab. They do not
have that property uniformly as its lower time tends to zero.

For completeness, the mechanism is the preceding head calculation.
On a spatial radius $L^2$ the clock geometry has normalized error
$O(LX^{-c t_0}+L^4/X)$; hence \eqref{cont:head-bounds} with
$J\asymp L^3$ gives the forcing floor. The complete inverse-square
exterior contributes $O(L^{-3})$ to the normalized interaction.
Inserting these quantities into \eqref{cont:barrier-bound} gives a
relative loss $O(\Delta/L)$ for fixed $\Delta$ within the supplied
slab. The localization theorem already guarantees reality and simplicity
in that exterior. Consequently this acquisition explains the dynamics
and the available rates, but does not extend the zero-free/collision-free
coverage beyond the imported theorem. The new question is precisely
how to acquire useful inputs after that localization estimate ceases
to provide them.
